\section{Finite group monoid semirings}
\label{sec:finite-group-projectives}

We prove Theorem~\ref{thm:finite-group-projectives}. Put $S=R_0$ and
$T=S[\Gamma]$. A left $T$-semimodule is exactly an $S$-semimodule
with a $\Gamma$-action by $S$-linear automorphisms. Restriction gives
one direction; in the other, set
\[
 \left(\sum_{g\in\Gamma}a_g g\right)m
   =\sum_{g\in\Gamma}a_g(gm).
\]
The old zero coefficient remains an occupied coefficient, so its action
is a fiber zero rather than an omitted summand. The distributive and
unit laws follow from those for $S$ and the commuting group action.

The support case is related to the study of group representations over
idempotent semifields. Jun, Mincheva and Tolliver
\cite[arXiv version~3, \S1.1 and Definition~3.1]{JMT} classify actions on free underlying
semifield modules; their Corollary~4.9 and Lemma~4.11 give further
Boolean-module results. Our modules need not be free over $\B$ or
$R_0$. Macpherson's $\mathbb F_1$-domain hypothesis
\cite[Definition~4.9, p.~14]{Macpherson} excludes nontrivial
multiplicative torsion. Stabilizer group rings also occur in the
classical splitting theory for modules over categories whose
endomorphisms are isomorphisms, recalled by Fiore, L\"uck and Sauer
\cite[pp.~306--307, splitting functors and Theorem~1.5]{FLS}. The proof here
constructs the required equivariant splittings directly over an
arbitrary unital coefficient ring.

\begin{proof}[Proof of Theorem~\ref{thm:finite-group-projectives}]
As an $S$-semimodule, $T$ is free on the finite set $\Gamma$.
Consequently a finite projective $T$-semimodule $M$ is finite projective
over $S$. By Lemma~\ref{lem:noncomm}, its support is $\Id(P)$ for a
finite poset $P$, and its group transitions are injective. Identify all
its groups with submodules of the top group. The classification gives
\[
 M_I+M_J=M_{I\cup J},\qquad M_I\cap M_J=M_{I\cap J}.
\]
The group action induces order automorphisms of $P$. Write
\[
 I_p=\downarrow p,\qquad I_p^-=I_p\setminus\{p\},\qquad
 V_p=M_{I_p}/M_{I_p^-}.
\]
This quotient is an $R[\Gamma_p]$-module. Its underlying $R$-module
is finite projective; we must prove the stronger stabilizer assertion.

For one $p$ in each orbit choose a finite free $R[\Gamma_p]$-module
$W_p$ surjecting onto $M_{I_p}$. Such a cover exists because $M_{I_p}$
is finitely generated over $R$: send a group-ring basis to an
$R$-generating family and extend equivariantly. Transport these modules
and maps along their orbits. The diagram
\[
 N_I=\bigoplus_{q\in I}W_q
\]
with its transported action is a $T$-semimodule with the same support.
The chosen maps and the group transitions define a $T$-map $f:N\to M$
inducing the identity on support. It is surjective on every group:
$M_I$ is the sum of the principal groups $M_{I_q}$ for $q\in I$,
and each of those is covered. Projectivity of $M$ supplies a section
with identity support. Passing to the principal birth quotients gives
\[
 V_p\longrightarrow W_p\longrightarrow V_p
\]
with identity composite, linearly over $R[\Gamma_p]$. Thus $V_p$ is
finite projective over that group ring. This argument does not require
$N$ to be projective.

Choose an $R[\Gamma_p]$-linear section of $M_{I_p}\to V_p$ for each
orbit representative, and transport it along $\Gamma$. Stabilizer
linearity makes the transported section independent of the chosen
coset representative. The sum of these sections identifies
$\bigoplus_{p\in I}V_p$ with $M_I$. Indeed, for a maximal $p\in I$
and $J=I\setminus\{p\}$, the two displayed identities give
\[
 M_I=M_J+M_{I_p},\qquad
 M_J\cap M_{I_p}=M_{I_p^-}.
\]
The selected section therefore yields $M_I=M_J\oplus V_p$, and finite
induction proves the claim. The identifications commute with all
transitions and with $\Gamma$. This proves necessity and the asserted
form of the module.

For sufficiency, start with the stated data. For each orbit representative
$p_a$, let $H_a=\Gamma_{p_a}$ and choose a split pair
\[
 V_{p_a}\xrightarrow{i_a}R[H_a]^{m_a}
 \xrightarrow{r_a}V_{p_a},\qquad r_ai_a=1,
\]
with $m_a\ge1$. For a zero label take $m_a=1$ and zero maps; the
coordinate still records the support point. Let $n=\sum_a m_a$.
The underlying $S$-free coordinates of $T^n$ are $(g,a,j)$ for
$g\in\Gamma$ and $1\le j\le m_a$. Define support maps by
\[
 \sigma(I)=\{(g,a,j):gp_a\in I\},\qquad
 \rho(U)=\bigcup_{(g,a,j)\in U}\downarrow(gp_a).
\]
They preserve unions and $\Gamma$-actions, hence are
$\mathcal P(\Gamma)$-linear, and $\rho\sigma=1$.

Induce the split pairs from $H_a$ to $\Gamma$. The resulting top group
maps are a split pair between
\[
 \bigoplus_a R[\Gamma]\otimes_{R[H_a]}V_{p_a}
 \quad\text{and}\quad
 \bigoplus_a R[\Gamma]^{m_a}.
\]
The first module is the direct sum of all the transported labels.
The section sends the label at $gp_a$ into coordinates $(gh,a,j)$
with $h\in H_a$; all have attached point $ghp_a=gp_a$.
It therefore takes the group over $I$ into the free group over
$\sigma(I)$. The retraction sends coordinate $(g,a,j)$ into the label
at $gp_a$, and hence into the group over $\downarrow(gp_a)$.
Both maps respect transitions and are $R$-linear and
$\Gamma$-equivariant. Together with $\sigma,\rho$, they give
$T$-linear maps
\[
 M\longrightarrow T^n\longrightarrow M
\]
whose composite is the identity on groups and supports. Thus $M$ is
finite projective. If $P$ is empty, the zero-fold free module suffices.

Finally any $T$-isomorphism induces a $\Gamma$-equivariant isomorphism
of the support posets and of all the intrinsic quotients $V_p$ with
their stabilizer actions. Conversely such data identify the split
diagrams and thus the $T$-semimodules. This proves precisely the
claimed isomorphism-class classification; it does not assert that
all homomorphisms are diagonal on the labels.
\end{proof}

For example, take $R=\Z$, $\Gamma\ne1$ finite, and $P$ a single
fixed point. The trivial label $\Z$ fails the stabilizer condition.
Indeed an equivariant section of the augmentation
$\Z[\Gamma]\to\Z$ would send $1$ to an invariant element
$c\sum_{g\in\Gamma}g$, whose augmentation is $c|\Gamma|\ne1$.
The regular label $\Z[\Gamma]$ always satisfies the condition.
When $R=0$, all labels vanish and the explicit support retraction
above proves the asserted classification over $\B[\Gamma]$.

\begin{corollary}\label{cor:finite-group-kzero}
Let $\mathcal C_{R,\Gamma}$ be the isomorphism classes of the nonempty
data in Theorem~\ref{thm:finite-group-projectives} for which $\Gamma$
acts transitively on the connected components of the comparability graph.
Then
\[
 V(R_0[\Gamma])\simeq\N^{(\mathcal C_{R,\Gamma})},\qquad
 K_0(R_0[\Gamma])\simeq\Z^{(\mathcal C_{R,\Gamma})}.
\]
The labels retain their actual projective $R[\Gamma_p]$-module
isomorphism classes. Every finitely generated strongly projective
$R_0[\Gamma]$-semimodule is free, and
\[
 SV(R_0[\Gamma])\simeq\N,\qquad SK_0(R_0[\Gamma])\simeq\Z.
\]
These assertions include $R=0$.
\end{corollary}
\begin{proof}
Direct sum is disjoint union of the decorated $\Gamma$-posets.
The comparability graph has a unique set of connected components,
which $\Gamma$ permutes. The unions belonging to individual
$\Gamma$-orbits of components give a unique decomposition into
nonempty $\Gamma$-invariant pieces. Such a piece admits no further
disjoint invariant decomposition. The classification therefore gives
the free commutative monoid, and its group completion is the displayed
free abelian group. No stabilization of individual ring labels is made.

The decorated poset of $T^n$ is the disjoint union of $n$ regular
$\Gamma$-orbits, with the discrete order, trivial stabilizers, and
label $R$. If $M\oplus N\simeq T^n$, the poset of $M$ is an invariant
subunion of these regular orbits with their unchanged labels. It is
therefore the poset of $T^k$ for some $k$. Conversely free modules
are strongly projective. Distinct free ranks have distinct support
cardinalities $2^{|\Gamma|k}$, also when the ring is zero. This proves
the strong-projective formulas.
\end{proof}
